\documentclass[11pt,a4paper]{article}
\usepackage[T1]{fontenc}
\usepackage[utf8]{inputenc}
\usepackage[french]{babel}
\usepackage{lmodern}
\usepackage[margin=2.2cm]{geometry}
\usepackage{amssymb}
\usepackage{xcolor}
\usepackage[most]{tcolorbox}
\usepackage{booktabs,tabularx,array}
\usepackage{enumitem}
\usepackage{tikz}
\usetikzlibrary{positioning,arrows.meta}
\usepackage[colorlinks=true,linkcolor=blue!50!black,urlcolor=blue!50!black]{hyperref}

\setlength{\parindent}{0pt}
\setlength{\parskip}{0.55em}
\setlist{itemsep=0.2em,topsep=0.25em}
\newcolumntype{L}{>{\raggedright\arraybackslash}X}

\newcommand{\cmd}[1]{\texttt{#1}}
\newcommand{\proj}{\textsf{Incognita}}
\newcommand{\lmlx}{\textsf{LMLExplore}}

% Fiche d'une tâche
\newtcolorbox{tache}[1]{enhanced,breakable,colback=blue!3,colframe=blue!55!black,
  fonttitle=\sffamily\bfseries,title={#1}}
\newtcolorbox{manque}{enhanced,breakable,colback=red!3,colframe=red!55!black,
  fonttitle=\sffamily\bfseries\small,title={Ce qui manque}}
\newtcolorbox{note}{enhanced,breakable,colback=green!4,colframe=green!45!black,
  fonttitle=\sffamily\bfseries\small,title={À retenir}}

% Rubriques d'une fiche
\newcommand{\rub}[1]{\par\smallskip\textsf{\textbf{#1.}}~}

\title{Les tâches d'\proj{}\\[0.3em]
  \large Chaque tâche, ses critères, et la comparaison avec \lmlx{}}
\author{Adrien Prevost\\ \small rédigé avec l'aide de Claude (Anthropic)}
\date{13 novembre 2026 --- état du dépôt après les PR \#1 à \#4}

\begin{document}
\maketitle

\begin{note}
Ce document décrit ce que font les scripts et les consignes \emph{tels qu'ils sont écrits}
(\cmd{AGENTS.md}, \cmd{agents/roles/*.md}, \cmd{scripts/agent.sh}, \cmd{scripts/loop.sh}, la CI).
Les scripts d'agents ont tourné pour la première fois le 13/11 au soir (essai du bibliographe) ;
le reste n'a pas encore été observé en conditions réelles.
\end{note}

\tableofcontents

%======================================================================
\section{Vue d'ensemble}
%======================================================================

Il y a trois sortes d'acteurs :
\begin{itemize}
\item \textbf{toi} (les humains) : écrire et approuver les programmes, décider pendant la calibration,
  fusionner, auditer par échantillon ;
\item \textbf{les agents}, une session Claude par tâche, lancée par \cmd{scripts/agent.sh} ;
\item \textbf{les automates} : la boucle \cmd{scripts/loop.sh} (choix des tâches, quota, fusion
  automatique), la CI de GitHub (compilation, audit, PDF).
\end{itemize}

\subsection*{Le cycle d'un résultat}

\begin{center}
\begin{tikzpicture}[
  b/.style={draw,rounded corners,align=center,font=\small,minimum height=0.9cm,inner sep=4pt},
  a/.style={-{Stealth[length=2.2mm]},thick}]
  \node[b,fill=gray!10] (p) {Programme\\approuvé (toi)};
  \node[b,fill=blue!8,right=0.5cm of p] (q) {Question\\proposée};
  \node[b,fill=blue!8,right=0.5cm of q] (v) {Question validée\\énoncé gelé};
  \node[b,fill=orange!10,right=0.5cm of v] (r) {Recherche\\preuve / réfutation};
  \node[b,fill=blue!8,below=0.9cm of r] (pr) {PR relue\\CI verte};
  \node[b,fill=green!12,left=0.5cm of pr] (f) {Fusion\\(toi ou auto)};
  \node[b,fill=green!6,left=0.5cm of f] (bib) {Audit\\bibliographe};
  \node[b,fill=green!6,left=0.5cm of bib] (j) {Jardinage\\confiance};
  \draw[a] (p)--(q); \draw[a] (q)--(v); \draw[a] (v)--(r); \draw[a] (r)--(pr);
  \draw[a] (pr)--(f); \draw[a] (f)--(bib); \draw[a] (bib)--(j);
\end{tikzpicture}
\end{center}

Les résultats \emph{connus} dont un programme a besoin suivent un chemin plus court : formalisation
directe (énoncé gelé tiré de la source), relue pour sa fidélité à la source, sans question.

\subsection*{Ordre de priorité de la boucle}

À chaque tour, \cmd{scripts/loop.sh} : (0) fusionne ce qui peut l'être automatiquement, (0') attend
que le quota le permette, puis lance \emph{une} tâche, la première disponible dans cet ordre :

\begin{center}
\begin{tabularx}{\linewidth}{@{}rlL@{}}
\toprule
& Tâche & Quand \\
\midrule
1 & relire une PR & une PR ouverte n'a pas d'avis \cmd{Verdict:} sur son dernier commit \\
2 & jardiner & un tour sur dix \\
3 & valider une question & une question \cmd{proposed} n'a pas encore d'avis du relecteur \\
4 & bibliographe (audit) & un résultat du catalogue a \cmd{novelty = "unchecked"} \\
5 & bibliographe (routine) & un tour sur quatre \\
6 & recherche & une question \cmd{open} existe, un tour sur deux \\
7 & formaliser un résultat connu & deux tours sur trois parmi les restants \\
8 & proposer une question & sinon \\
\bottomrule
\end{tabularx}
\end{center}

Les tâches 2 à 8 ne portent que sur les programmes \textbf{approuvés} (aujourd'hui : BatchedBandits).

%======================================================================
\section{Les tâches des agents}
%======================================================================

Chaque fiche donne : le \emph{rôle} (fichier de consignes), le \emph{déclencheur}, ce que l'agent
\emph{lit}, ce qu'il \emph{écrit}, les \emph{critères} qu'il doit appliquer, et les \emph{contrôles}
qui ne dépendent pas de lui (script, CI, relecteur). Tous les agents peuvent lire le web (recherche et
pages), mais n'ont ni réseau en ligne de commande, ni \cmd{git push}, ni \cmd{gh} : c'est le script qui
pousse la branche et ouvre la PR. Tous ceux qui écrivent terminent par un résumé
(\cmd{.agents/summary.md}) qui devient la description de la PR et le PDF.

%----------------------------------------------------------------------
\begin{tache}{1. Proposer une question --- rôle \cmd{propose} (modèle auteur)}
\rub{Déclencheur} rien de plus prioritaire ; programme approuvé tiré au hasard.
\rub{Lit} le programme et ses fichiers compagnons, l'index et toutes les questions existantes, le
catalogue.
\rub{Écrit} un fichier \cmd{programmes/<P>/questions/Q<nnn>-<titre>.md} (gabarit
\cmd{templates/QUESTION.md}), statut \cmd{proposed}, << décision du propriétaire requise : oui >>, et une
ligne dans l'index. Éventuellement un brouillon dans \cmd{Exploration/}.
\rub{Critères}
\begin{itemize}
\item entrer dans une direction ouverte du programme et satisfaire \emph{chacun} de ses critères de
  validation (vérifiés par l'agent lui-même avant d'écrire) ;
\item ne dupliquer aucune question existante, ne pas être un corollaire du catalogue (essayer d'abord
  une preuve courte ; si elle marche, choisir autre chose) ;
\item chercher la réponse dans la littérature ; si elle est connue, changer de question ou la
  transformer en résultat connu à formaliser ;
\item préférer ce qui compterait pour le programme à ce qui est seulement faisable ; une question
  dont la réponse probable est << non >> est acceptable ;
\item contenir un exemple qui satisfait toutes les hypothèses à la fois, et la variante de repli.
\end{itemize}
\rub{Contrôles} le script refuse tout fichier hors des questions et d'\cmd{Exploration/} ; la question
passe ensuite par la tâche 2.
\end{tache}

%----------------------------------------------------------------------
\begin{tache}{2. Valider une question --- rôle \cmd{review-question} (modèle relecteur)}
\rub{Déclencheur} une question \cmd{proposed} sans avis du relecteur (priorité haute).
\rub{Lit} uniquement la question, son programme et le catalogue ; pas l'historique git ni ce que le
proposeur a écrit ailleurs (indépendance).
\rub{Écrit} son verdict (\cmd{accept} ou \cmd{reject}, modèle, date) et un rapport critère par
critère dans la question ; s'il accepte, l'\textbf{énoncé gelé} en Lean
(\cmd{Statements/<P>/Q<nnn>.lean}, avec \cmd{sorry}), vérifié par la compilation.
\rub{Critères}
\begin{itemize}
\item chaque critère de validation du programme : << rempli >> ou << non rempli >>, avec une phrase ;
\item tentative de preuve courte à partir du catalogue et des dépendances : si elle réussit, rejet
  (corollaire) ;
\item recherche de la réponse dans la littérature : si elle est connue, rejet avec la référence ;
\item intérêt pour le programme (section << ce qui serait un résultat important >>) ;
\item l'énoncé gelé n'utilise que des définitions de la bibliothèque et des dépendances.
\end{itemize}
\rub{Contrôles} pendant la \textbf{calibration}, le statut reste \cmd{proposed} : c'est toi qui décides,
et l'accord entre ton avis et celui du relecteur est compté. L'ajout d'un fichier dans
\cmd{Statements/} n'est jamais fusionné automatiquement.
\end{tache}

%----------------------------------------------------------------------
\begin{tache}{3. Recherche sur une question ouverte --- rôle \cmd{research} (modèle auteur)}
\rub{Déclencheur} une question \cmd{open} existe (un tour sur deux).
\rub{Lit} la question et ses impasses, le programme, l'énoncé gelé, le catalogue.
\rub{Écrit} selon l'issue : brouillons dans \cmd{Exploration/} ; preuve dans la bibliothèque sous le nom
\cmd{Incognita.<P>.<nom>} ; rédaction lisible avec rubrique << Attribution >> ; entrée du catalogue
avec \cmd{novelty = "unchecked"} ; statut \cmd{proved}, \cmd{refuted} ou \cmd{variant-proved} ; sinon une
ligne d'impasse par approche ratée.
\rub{Critères}
\begin{itemize}
\item prouver \emph{exactement} l'énoncé gelé, jamais l'affaiblir ;
\item une réfutation (preuve de la négation) est un résultat à part entière ;
\item une variante exige une raison formelle : contre-exemple à l'original, ou preuve que l'hypothèse
  ajoutée est nécessaire, ou étiquette << original toujours ouvert >> ;
\item bibliothèque : aucun \cmd{sorry}, axiomes standard seulement, aucun avertissement ;
\item toute nouvelle définition : témoin non trivial, retour aux cas connus, non-vacuité conjointe
  des hypothèses, finesse quand c'est possible ;
\item ne jamais déclarer le résultat nouveau soi-même ; citer seulement des références lues.
\end{itemize}
\rub{Contrôles} CI : compilation sans \cmd{sorry} ni avertissement, \cmd{lake exe audit} (axiomes, et
preuve du même type que l'énoncé gelé, testé : il rejette les deux fautes) ; relecture (tâche 6) ;
audit du bibliographe après fusion.
\end{tache}

%----------------------------------------------------------------------
\begin{tache}{4. Formaliser un résultat connu --- rôle \cmd{formalise} (modèle auteur)}
\rub{Déclencheur} un programme approuvé a une liste << Known results to formalise >> (deux tours sur
trois parmi les restants).
\rub{Lit} le programme, le catalogue, la fiche de lecture de la source.
\rub{Écrit} l'énoncé gelé avec la citation exacte \cmd{[CLÉ, Théorème X]} et les différences de forme ;
la preuve dans la bibliothèque ; l'entrée du catalogue (\cmd{source}, \cmd{novelty = "known: ..."}) ;
au besoin la fiche de lecture.
\rub{Critères}
\begin{itemize}
\item prendre le premier élément non fait, non marqué \emph{on hold}, absent des dépendances ;
  s'il est trop gros, une partie autonome (une définition et son API, un lemme) ;
\item formaliser à partir de la source, jamais de mémoire : l'énoncé lu doit être recopié dans
  \cmd{literature/notes/} ;
\item suivre les noms et formes d'énoncés prévus par les feuilles de route de \lmlx{}, pour qu'une
  adoption ultérieure soit une suppression plus un import ;
\item si la source n'est pas publiée et que la preuve échoue à cause de l'énoncé : le signaler dans une
  question, ne jamais affaiblir en silence.
\end{itemize}
\rub{Contrôles} les mêmes que la tâche 3 ; la relecture vérifie la fidélité à la source plutôt que la
nouveauté.
\end{tache}

%----------------------------------------------------------------------
\begin{tache}{5. Bibliographe --- rôle \cmd{bibliographer} (modèle auteur)}
\rub{Déclencheur} un résultat non audité (priorité haute), sinon un tour sur quatre.
\rub{Fait} une seule de ces trois choses, dans cet ordre :
\begin{enumerate}
\item \textbf{audit d'attribution} d'un résultat fusionné : le résultat est-il connu, cas particulier
  d'un résultat connu ? d'où vient chaque technique de la preuve (la plus ancienne source, en suivant les
  références dans les deux sens) ? Il fixe \cmd{novelty} à \cmd{new}, \cmd{rediscovery of [...]} ou
  \cmd{known: [...]}, complète la rubrique << Attribution >> et ajoute \cmd{attribution-audit} aux
  vérifications ;
\item \textbf{vérifier une référence} : titre, auteurs, revue, année, identifiant comparés à la page de
  l'éditeur ou d'arXiv ; correction de \cmd{bibliography.json}, passage à \cmd{checked} ;
\item \textbf{lire un passage cité} : recopier dans la fiche l'énoncé exact (hypothèses, conclusion,
  constantes), marquer le passage comme lu.
\end{enumerate}
\rub{Critères} ne consigner que ce qui a été lu dans la source ; ne rien remplir de mémoire ; ne pas
faire confiance aux rubriques << Attribution >> des autres, les vérifier ; citations courtes seulement.
\rub{Contrôles} le script refuse toute écriture hors de \cmd{literature/}, \cmd{catalogue/} et des
rédactions de questions.
\end{tache}

%----------------------------------------------------------------------
\begin{tache}{6. Relire une PR --- rôle \cmd{review-pr} (modèle relecteur)}
\rub{Déclencheur} une PR ouverte sans avis sur son dernier commit (priorité la plus haute).
\rub{Lit} le diff et la description de la PR, le code de la PR ; pas le raisonnement de l'auteur.
\rub{Écrit} rien dans le dépôt : un commentaire sur la PR, dont la première ligne est exactement
\cmd{Verdict: approve}, \cmd{Verdict: request\_changes} ou \cmd{Verdict: block}, puis les défauts du plus
grave au moins grave, avec ce qui les corrigerait.
\rub{Critères}
\begin{enumerate}
\item \textbf{rétrotraduction} : réécrire en mathématiques chaque énoncé Lean à partir du Lean seul, et
  comparer au texte (question, source, rédaction) ;
\item \textbf{lecture adversariale} : énoncé vide, quantificateurs, contenu caché dans les hypothèses,
  cas dégénérés (ensemble vide, zéro, division par zéro, soustraction dans $\mathbb N$, indices à partir
  de 0) ;
\item \textbf{règles d'\cmd{AGENTS.md}} : énoncés gelés intacts, variantes justifiées, définitions
  contrôlées, pas de définition évaluative, aucun théorème perdu, catalogue et question à jour, chemins
  humains intacts ;
\item \textbf{réutilisation} : rien de redéfini qui existe déjà ;
\item \textbf{citations} : chaque \cmd{[CLÉ, ...]} existe, passage lu ou marqué << to be checked >>,
  rubrique Attribution présente, rien d'annoncé comme nouveau avant l'audit.
\end{enumerate}
\rub{Contrôles} le script n'accepte qu'un avis qui commence par \cmd{Verdict:} ; le relecteur tourne sur
un autre modèle que l'auteur par défaut (Sonnet contre Opus).
\end{tache}

%----------------------------------------------------------------------
\begin{tache}{7. Jardiner --- rôle \cmd{garden} (modèle auteur)}
\rub{Déclencheur} un tour sur dix.
\rub{Fait} une consolidation, la plus utile : une implication entre résultats du catalogue (cas
particulier devenu corollaire, l'énoncé restant), une borne inférieure et une borne supérieure de même
classe confrontées (si Lean prouve \cmd{False}, une définition est fausse : le signaler), un doublon avec
les dépendances supprimé, le catalogue mis à jour.
\rub{Critères} englobement seulement à contenu quantitatif au moins aussi bon ; \textbf{aucun théorème
perdu} (l'ancien énoncé doit rester prouvable) ; ne rien changer si rien d'utile.
\rub{Contrôles} CI et audit ; relecture ; seul un humain peut supprimer un énoncé.
\end{tache}

%======================================================================
\section{Tes tâches et celles des automates}
%======================================================================

\begin{center}
\begin{tabularx}{\linewidth}{@{}lL@{}}
\toprule
Qui & Quoi \\
\midrule
Toi & écrire, relire et approuver les programmes ; décider des questions pendant la calibration ;
  fusionner les PR (avec le PDF résumé) ; auditer de temps en temps quelques questions acceptées et
  quelques fiches de lecture ; décider des suppressions d'énoncés ; protéger \cmd{main}. \\
\addlinespace
Boucle & choisir la tâche (section 1) ; attendre le quota avec le calcul du worker Tau Ceti
  (\cmd{INCOGNITA\_PACE}) ; fusionner automatiquement les catégories autorisées
  (\cmd{INCOGNITA\_AUTOMERGE}, vide par défaut ; jamais \cmd{Statements/} ni les chemins humains). \\
\addlinespace
Script d'agent & copie de travail isolée ; outils restreints ; refus des chemins interdits ; push et
  ouverture de la PR ; publication de l'avis du relecteur. \\
\addlinespace
CI & compilation de la bibliothèque sans \cmd{sorry} ni avertissement, des énoncés gelés et des
  brouillons ; \cmd{lake exe audit} ; linters de Mathlib ; versions épinglées comme \lmlx{} ; racines
  complètes ; interdiction d'importer \cmd{Exploration} ; PDF résumé de chaque PR. \\
\bottomrule
\end{tabularx}
\end{center}

%======================================================================
\section{Comparaison avec \lmlx{}}
%======================================================================

\subsection{Ce que fait le worker Tau Ceti}

D'après les consignes de \lmlx{} et ton guide du worker, chaque ronde fait \emph{une} tâche, la plus
utile disponible. Il y en a \textbf{trois}, pas deux :
\begin{enumerate}
\item \textbf{s'occuper de ses propres PR} : les rebaser, réparer la CI, répondre aux relecteurs (c'est
  la priorité) ;
\item \textbf{relire les PR des autres}, selon des grilles fixes : correctness, reuse, scope,
  attribution, API design, generality, placement, naming, documentation, proof quality. Chaque grille
  donne un verdict \cmd{approve}, \cmd{request\_changes} ou \cmd{block} ;
\item \textbf{écrire une nouvelle PR} qui fait avancer une cible d'une feuille de route.
\end{enumerate}
Autour : les humains écrivent et relisent les feuilles de route ; une PR est fusionnée automatiquement
quand toutes les grilles approuvent et que la CI est verte ; le worker d'un mainteneur peut finir une PR
inactive depuis un jour ; un script signale les doublons avec les dépendances.

\subsection{Tableau comparatif}

\begin{center}
\small
\begin{tabularx}{\linewidth}{@{}lLL@{}}
\toprule
& \lmlx{} & \proj{} \\
\midrule
Origine des énoncés & feuilles de route humaines, résultats publiés & programmes humains, questions proposées par des agents \\
Vérité de l'énoncé & supposée (vient de la littérature) & inconnue : preuve, réfutation ou variante justifiée \\
Validation de ce qu'on fait & revue humaine des feuilles de route & programme approuvé par toi, question validée par un agent (toi pendant la calibration) \\
Tâches des agents & 3 : corriger ses PR, relire, écrire une PR & 7 : proposer, valider une question, chercher, formaliser, bibliographe, relire, jardiner \\
Relecture & 10 grilles, une par aspect & 1 relecteur qui couvre 5 aspects (dont citations) \\
Indépendance & mêmes modèles pour écrire et relire & relecteur sur un autre modèle par défaut, ne voit pas le raisonnement \\
Énoncé gelé & non (fidélité jugée par la grille correctness) & oui, vérifié par \cmd{lake exe audit} \\
Résultats négatifs & pas de place prévue & réfutations et impasses consignées \\
Bibliographie & citations dans les docstrings, liens vers la littérature & bibliographie vérifiée, passages lus, audit d'attribution \\
Consolidation & doublons avec les dépendances & jardinage : englobement, aucun théorème perdu \\
Confiance & fusionné ou non & niveaux : fusionné, contrôlé, utilisé, établi \\
Fusion & automatique (toutes les grilles + CI) & par toi ; automatique par catégorie si tu l'actives \\
Quota & \cmd{--pace} du worker & le même calcul, via \cmd{tauceti status} \\
Résumé pour l'humain & description de la PR & description + PDF avec les énoncés Lean \\
\bottomrule
\end{tabularx}
\end{center}

\begin{manque}
\begin{itemize}
\item \textbf{<< Corriger ses propres PR >> n'existe pas encore comme tâche.} \cmd{AGENTS.md} la met en
  première priorité, mais aucun rôle ni la boucle ne l'implémentent : une PR qui reçoit
  \cmd{request\_changes} ou dont la CI échoue reste bloquée jusqu'à ce que toi ou moi intervenions. C'est
  la première chose à ajouter (rôle \cmd{fix-pr} : rebaser, réparer la CI, traiter les défauts signalés,
  puis nouvelle relecture).
\item \textbf{Un seul relecteur} contre dix grilles chez \lmlx{} : moins coûteux, mais chaque aspect est
  regardé moins attentivement. On pourra découper la relecture en grilles si la calibration montre des
  oublis.
\item \textbf{Double formalisation} des énoncés principaux et \textbf{niveaux de confiance} : écrits dans
  les règles, mais aucune tâche ne les déclenche encore.
\end{itemize}
\end{manque}

%======================================================================
\section{Premier essai réel (13 novembre 2026)}
%======================================================================

Sur le serveur, une tâche de bibliographe puis une relecture, lancées à la main avec
\cmd{scripts/agent.sh}.

\begin{itemize}
\item \textbf{Bibliographe} (Opus, quelques minutes). Il a constaté qu'aucun résultat n'était à
  auditer, puis a vérifié la première référence non vérifiée citée par un programme, \cmd{AG12}
  (Agrawal--Goyal, COLT 2012) : titre aligné sur l'éditeur, nom exact des actes, volume, pages,
  éditeurs, identifiant arXiv. Il a ouvert la PR \#5 avec un résumé correct, et a signalé de lui-même
  une ambiguïté : le programme RegretShape cite << Agrawal--Goyal >> sans année, alors que les mêmes
  auteurs ont des articles de 2013 et 2017 plus proches du sujet (termes d'ordre 2).
\item \textbf{Relecteur} (Sonnet, sans voir le raisonnement de l'auteur). Avis
  \cmd{Verdict: request\_changes}, publié sur la PR, avec quatre défauts dont deux sérieux : le statut
  \cmd{checked} ne s'appuie sur aucune trace de la comparaison avec la source, et l'identification de
  la référence n'est qu'une inférence. Il a aussi dit honnêtement qu'il n'avait pas ouvert les pages
  sources lui-même.
\item \textbf{Ce qui a fonctionné} : copie de travail, outils restreints et accès web, commit avec
  l'identité des agents, contrôle des chemins, push, PR, résumé, PDF en CI, avis du relecteur.
\item \textbf{Bug trouvé} : le script tentait de supprimer la branche de relecture avant la copie de
  travail qui l'utilisait (sans conséquence, mais une branche restait sur le serveur) ; corrigé.
\item \textbf{Ce que l'essai montre} : la PR \#5 attend maintenant des corrections, et aucune tâche ne
  les fera, puisque << corriger ses propres PR >> n'existe pas encore (section 4).
\end{itemize}

\end{document}
